\documentclass[12pt,onecolumn]{article}

\usepackage[utf8]{inputenc}
\usepackage[T1]{fontenc}
\usepackage{lmodern}
\usepackage[margin=1in]{geometry}
\usepackage{setspace}
\usepackage{microtype}
\usepackage{amsmath,amssymb}
\usepackage[numbers,sort&compress]{natbib}
\usepackage[colorlinks=true,linkcolor=blue,citecolor=blue,urlcolor=blue]{hyperref}
\usepackage{orcidlink}
\usepackage{fancyhdr}

\setstretch{1.10}
\setlength{\headheight}{15pt}

\pagestyle{fancy}
\fancyhf{}
\lhead{Space Is Not a Box}
\rhead{John C. W. McKinley}
\cfoot{\thepage}

\title{
\textbf{Space Is Not a Box}\\[0.4em]
\large Timeless Law, Universal Actuality,
and the Registration of Space and Motion
}

\author{
John C. W. McKinley\,\orcidlink{0009-0005-7097-5035}\\
\small Independent Researcher
}

\date{October 2026}

\begin{document}
\maketitle

\noindent\footnotesize\textit{Working draft published at \href{https://johncwmckinley.com}{johncwmckinley.com}. Not deposited on Zenodo; subject to revision.}\normalsize

\begin{abstract}
The ordinary understanding of space begins with a container. Physical objects occupy locations within it, move between those locations, and leave histories that observers remember. Although special and general relativity replace Newtonian absolute space with a more sophisticated geometric description, the underlying mental picture of objects traveling through a completed spatial arena often remains.

This paper develops an alternative interpretation within the Timeless Light Model (TLM). The proposed foundation of the model is the totality of timeless governing laws, not space, time, motion, or observation. These laws are not themselves spatial or temporal constituents. They govern the admissible relationships, interactions, and outcomes of physical reality. The manifestation of space in a human observer or recording device is one application of this broader lawful architecture.

The account distinguishes physical actuality, registration, and memory. Physical actuality consists of the lawful physical conditions and interactions that are actual, independently of whether a particular observer knows about them. Registration is the causally constrained physical reveal through which consequences become expressed in interactions and available to physical systems. Records preserve information concerning earlier registrations. In mechanical and electronic devices, such records are identifiable physical configurations. Human recollection likewise provides present access to information concerning earlier experiences, without requiring this paper to determine the ultimate physical or nonphysical nature of subjective memory. Inspection of records permits the reconstruction of distance, motion, duration, and physical history.

A meteor observed crossing the sky illustrates the interpretation. The observer's visual system undergoes a sequence of lawful physical changes associated with the meteor, its environment, and the arrival of electromagnetic information. The preserved consequences of those changes support the experience and later reconstruction of a meteor moving through space. That reconstruction is physically constrained and predictive. It does not, by itself, require an independent spatial container or a substance traveling through such a container.

The paper further adopts the TLM proposition that physical actuality is singular: there is one universal present, while elapsed proper time is local and information becomes available through relativistically constrained registration. This universal-present proposition is ontological and is not identified with an observer-independent global simultaneity coordinate.

The result is a standalone foundations argument. It does not modify established equations or predictions of quantum mechanics, special relativity, or general relativity. Its contribution is to distinguish timeless law from physical registration, and registered spatial experience from an independently existing spatial container, while identifying the additional mathematical work required to turn this interpretation into a predictive physical theory.
\end{abstract}

\section{Introduction: Watching a Meteor}

Consider a person named Bob standing outdoors on a clear night. A meteor becomes visible overhead, crosses part of the sky, and disappears. Bob describes the experience in familiar terms: a physical object moved through space.

Nothing is wrong with this description as an account of the observable phenomenon. The apparent trajectory can be photographed, its changing angular position can be measured, and its physical interactions with the atmosphere can be analyzed. Other observers may contribute additional measurements that permit the reconstruction of its path, speed, altitude, and possible impact location.

The question addressed here is not whether these measurements work. It is what their success requires us to believe exists.

The ordinary picture places the meteor inside a completed spatial container. The meteor occupies one position, then another, then another. Its motion is imagined as the continuous displacement of an independently existing object through an independently existing arena. Bob watches the object travel, and his memory preserves an account of where it has been.

This picture brings together several distinct ideas without necessarily distinguishing them. It assumes the existence of a spatial container, the persistence of a traveler-substance, the passage of time, the occurrence of physical interactions, and the retention of a history. It then treats the observer's experience of motion as evidence for the entire picture.

Modern physics does not require all of these assumptions. Special relativity denies a preferred inertial state of rest and makes the assignment of distant simultaneity dependent on the reference frame. General relativity describes spacetime through a dynamical metric rather than a fixed Newtonian arena. Quantum theory provides rules for physical outcomes that cannot generally be reduced to an inventory of classical objects with determinate properties independent of experimental conditions.

The rejection of an independently existing spatial container has important historical precedents. In his 1952 note to the fifteenth edition of \emph{Relativity: The Special and the General Theory}, Einstein wrote, ``Physical objects are not in space, but these objects are spatially extended'' \cite{EinsteinRelativity}. Carlo Rovelli subsequently stated, ``The physical world does not exist `in space','' describing spatial order in terms of relationships among physical systems \cite{Rovelli2004}.

The present paper accepts the importance of these relational insights but proposes a different explanatory emphasis. Within the Timeless Light Model, the totality of timeless governing laws is fundamental. The registration of spatial relationships is one application of that lawful structure. The laws do not reside within the geometry whose physical relationships they govern.

These developments motivate a reconsideration of the spatial-container picture, but they do not by themselves determine a unique ontology.

The Timeless Light Model proposes an interpretation in which the more fundamental starting point is lawful structure. The governing laws are not imagined as objects occupying space. They govern physical relationships and interactions. Those interactions produce consequences, some of which become registered in physical systems. Present physical records can then support the reconstruction of distance, motion, and history.

Bob's experience of the meteor is therefore not the foundation of the model. It is one use case through which the model's proposed architecture can be understood.

The central question is whether the lawful interactions, their geometric relationships, and their registered consequences are sufficient to account for what Bob observes, without adding an independently existing spatial box.

\section{The Totality of Timeless Law}

The foundation proposed by TLM is the totality of timeless governing laws.

This should not be understood as a single equation that somehow contains the entire universe, nor as a physical storage location for instructions. It refers to the complete lawful structure governing admissible physical relationships, interactions, and outcomes. Different situations bring different aspects of that structure into practical relevance.

The laws relevant to an atomic transition include quantum dynamics and the constraints governing the participating states and interactions. The laws relevant to planetary motion include gravitational dynamics and the appropriate relativistic geometry. The laws relevant to Bob's observation of a meteor include electromagnetic processes, atmospheric interactions, the causal propagation of radiation, and the physical organization of his sensory system.

These are not separate universes. They are particular applications of the proposed totality of lawful structure.

The term \emph{timeless} refers to the status attributed to the governing laws, not to the durations of the physical processes they govern. A law does not accumulate proper time, occupy a spatial position, or travel from one interaction to another. The applicability of the law governing an event is not a signal that must first arrive at the event.

This distinction does not authorize instantaneous physical communication. The laws may be timeless in the proposed ontology while the physical interactions they govern remain subject to relativistic causal restrictions. A law constraining information propagation does not itself constitute information propagating between physical systems.

A further distinction is needed between a law and the mathematical expression used to represent it. The equations of physics are human formulations that may be revised or improved. TLM's ontological proposal concerns the lawful structure those equations attempt to describe. It does not assume that the presently known equations exhaust that structure.

The claim that this structure is ontologically fundamental is therefore an interpretive premise, not a discovery established merely by writing down physical equations. The distinction matters because the proposal would otherwise confuse the mathematical description of a law with an independently demonstrated account of why that law exists.

The model begins with lawful structure rather than with an empty arena awaiting occupation. Space is not assumed in advance as the container in which the laws operate. Instead, the model asks how spatially describable relationships arise among the physical conditions and interactions governed by those laws.

The totality of timeless law is the model. The manifestation of space is one possible expression of its operation.

\section{The Laws Relevant to an Event}

Although the proposed lawful structure is comprehensive, a particular physical event is not described by indiscriminately invoking every known law. It is described through the laws and physical conditions relevant to that event.

For purposes of the present discussion, call these the \emph{lawful relationships under address}.

When Bob observes the meteor, the event under address is not limited to a point of light on his retina. It belongs to a connected physical situation that may include the meteor's atmospheric interactions, emitted radiation, causal propagation, absorption in Bob's eye, neural processing, and the formation of memory.

Each component of that situation is governed by applicable physical laws. The event under address identifies which lawful relationships are relevant to the question being investigated.

This provides a more precise meaning for the claim that space is relational.

The claim is not simply that an object has a position only relative to another object. Nor is it that every physical law should be relabeled a spatial relationship. Rather, spatial description concerns the geometric aspects of the lawful physical relationships relevant to events.

Those geometric aspects are expressed through such quantities as interval, separation, direction, curvature, and the relations governing trajectories. The laws governing the events are not themselves spatial objects or ingredients occupying the geometry.

The distinction is essential. The governing lawful structure and its spatial manifestation are not the same category of thing.

For example, the relationship between the arrival times of two signals at a detector may support a determination of direction or distance, subject to the appropriate assumptions. The signal arrivals are physical interactions. Their ordering is available through the detector's records. The geometric inference uses lawful relationships connecting the emission conditions, the propagation paths, and the detector.

At no point must the governing law itself be placed at a coordinate location.

The spatial description is an expression of the lawfully governed situation, not the home in which the laws reside.

\section{Physical Actuality Is Not Observation}

A registration-based interpretation can easily be misunderstood as the claim that observation creates physical reality. That is not the proposal advanced here.

Physical actuality and an observer's knowledge of physical actuality are different concepts.

The meteor interacts with the atmosphere whether or not Bob sees it. Its interactions may deposit energy, produce radiation, fragment material, and leave physical consequences. Bob's presence is not a condition for the governing laws to apply.

When radiation from the meteor reaches Bob, his visual system becomes another participant in the lawful physical situation. Bob is not a spectator standing outside reality. His eye and brain are physical systems participating in the same lawful structure. His remembered experience arises in association with those systems, but the argument does not require equating subjective memory with physical organization.

A mechanical detector may participate in an analogous way. A camera can record incoming radiation without consciousness. A collision can change the state of physical systems without producing a permanent or accessible record. The absence of a later observer does not make such interactions unreal.

TLM therefore distinguishes physical actuality from the accessibility of physical information.

What a detector registers depends on the lawful causal connections available to that detector. What physically exists is not limited to what the detector has already learned.

This distinction is particularly important in relativistic situations. An observer may receive information from a distant source only after a propagation interval. That delay restricts the observer's knowledge of distant conditions. It does not, by itself, establish that the distant source lacked a physical condition before the information arrived.

Delayed registration is not delayed existence.

The model's additional ontological proposition is that physical actuality is singular. The meaning and limitations of that proposition will be examined separately, because it cannot simply be identified with a universal clock or a preferred relativistic reference frame.

\section{Registration as the Lawful Reveal}

The word \emph{registration} serves an important role in this paper and must be used consistently.

In its most immediate physical sense, a registration is an interaction or actual outcome that produces a physical consequence. A photoreceptor changes state following absorption of radiation. A detector records a particle interaction. A collision transfers energy and momentum. These are lawful physical occurrences.

The broader expression \emph{relativistic registration} describes the causally constrained reveal through which such consequences become available within physical interactions and recording pathways. The reveal may involve a series of related actions, rather than a single measurement.

A physical record is not identical to the original registration. It is a condition that preserves some consequence of that registration and may itself arise through further interactions.

The distinction can be illustrated by a camera. Incoming radiation produces interactions in the sensor. Electronic processes convert those interactions into signals. The device may then preserve a pattern of information in memory. Later inspection of the stored data allows a person or another device to reconstruct aspects of the earlier physical situation.

Registration, record formation, and record inspection are distinguishable processes, even though they may be closely connected.

The term \emph{reveal} is appropriate because physical consequences become available to systems capable of interacting with them. But it should not imply that the lawful events are created by the later inspection of their records.

The reveal is governed by the causal relationships of special and general relativity. An event cannot ordinarily influence an arbitrarily distant detector without an appropriate causal connection. Radiation, matter, and other physical influences obey the relevant laws governing their propagation and interaction.

The invariant speed $c$ enters the causal structure of relativistic physics. It constrains the propagation of information and the geometry of the possible causal relationships between events.

TLM interprets the observable sequence of registered consequences through that causal structure. It does not remove the relativistic propagation requirements or replace them with an instantaneous transmission mechanism.

The governing law is not delayed by distance in the way a propagating signal is delayed. The physical reveal of information, however, remains subject to the relevant causal and geometric constraints.

This is one of the central distinctions of the model.

\section{What Bob Actually Experiences}

Return to Bob watching the meteor.

In the conventional interpretation, Bob sees a meteor traveling through a completed spatial arena. The positions are imagined to exist independently of his observation, and the meteor is imagined to move continuously between them.

The TLM interpretation examines the physical sequence more closely.

The meteor participates in lawful interactions with the atmosphere. Those interactions produce electromagnetic radiation and other physical consequences. Some of the radiation reaches Bob's retina. The resulting absorption events alter the states of photoreceptors, which contribute to further neural interactions. The visual system integrates information over finite intervals and may establish a persistent memory of the experience.

Bob's eye is not receiving a completed journey. It is undergoing lawful physical changes associated with the consequences reaching it.

As the meteor's visible conditions change, the lawful interactions contributing to Bob's sensory state also change. Different directions, intensities, colors, and apparent positions may become distinguishable through the visual system's responses.

The sequence is physical. It is not an invention by Bob.

Nevertheless, Bob does not directly inspect an independent spatial container. He possesses an evolving physical condition affected by a series of lawful interactions.

When the resulting records are compared, Bob can describe an ordered pattern. The meteor appeared in one direction, then in another, and then in another. The ordered pattern supports the description of an object moving across the sky.

This is what the ordinary expression \emph{the meteor moved through space} reports.

The experience and the reconstruction can be highly reliable. Multiple detectors may register compatible consequences and permit a common geometric reconstruction after accounting for their locations, orientations, and causal delays.

But the reliability of that reconstruction is not an independent proof of the spatial box imagined behind it.

The reconstruction establishes lawful physical relationships and their geometric description. Whether those relationships also require an independently existing spatial substance is a further ontological question.

Under TLM, the additional container is unnecessary.

\section{Present Records and Remembered Experience}

The meteor example depends on a feature of physical recording that is easy to overlook: device records are accessible through presently existing physical conditions. Human recollection is presently experienced, while its underlying ontology is not at issue here.

A photograph of the meteor does not preserve the original atmospheric event as an independently continuing occurrence. It preserves a physical record resulting from interactions associated with that event.

The same applies to an electronic recording. The present state of a memory circuit may encode a sequence of earlier measurements, but the stored data are not the earlier events themselves.

Human memory presents a different question. Bob presently remembers an earlier experience, and that recollection is associated with physical activity in his nervous system. Whether the subjective content of his memory is exhaustively reducible to a physical configuration is not established by this analysis, nor is such a conclusion necessary to the argument.

What matters here is that Bob's present recollection refers to an earlier event without requiring that event to continue existing as an independent physical actuality.

This is important because the interpretation of movement requires ordered relationships among records.

Bob may remember the meteor appearing near one star and then crossing toward another. A camera may preserve successive image frames. An instrument may retain time-tagged measurements. When these records are inspected, their contents can be arranged into a description of the observed history.

Such arrangements may be constrained by causal relationships, known physical laws, and the properties of the recording apparatus. They are not necessarily accurate merely because they exist. Human memories can be mistaken, instruments can fail, and limited records can support more than one possible history.

Physical reconstruction therefore remains subject to ordinary standards of empirical testing.

The ontological distinction is that present records do not require the continued physical existence of the earlier events they describe.

A worldline diagram offers a useful analogy. It may represent a sequence of events attributed to a physical system and permit the calculation of proper-time intervals. The mathematical representation is valuable without requiring every earlier portion of the represented history to persist as an independently existing physical trail.

Similarly, a sequence of meteor images may represent a history of lawful interactions without requiring all previous configurations of the meteor to continue existing alongside the present record.

The past can be represented without being retained as an independently completed physical world.

\section{There Is Only the Present}

TLM proposes that there is one physical actuality, which it calls the universal present.

The ordinary expression ``now, now, now'' is useful for conveying the intended interpretation. At each stage of a registered physical history, what is available is physical actuality. Earlier conditions may leave consequences that persist in present records. The earlier events themselves are not required to remain as separately existing physical configurations.

This does not mean that the universe possesses a master clock counting a sequence of universally synchronized instants.

The distinction is fundamental.

In special relativity, the simultaneity assigned to spatially separated events depends on the reference frame. For an inertial Lorentz transformation along one spatial direction,

\begin{equation}
t'=\gamma\left(t-\frac{vx}{c^2}\right),
\qquad
\gamma=\frac{1}{\sqrt{1-v^2/c^2}}.
\end{equation}

Two distant events that are simultaneous in one inertial coordinate system need not be simultaneous in another.

This is not merely a problem caused by observers receiving light at different times. It is a structural feature of relativistic spacetime.

The universal-present proposition therefore cannot be established by correcting observations for signal delay. Such corrections do not produce a preferred simultaneity relation accepted by all inertial observers.

TLM instead uses the universal present as an ontological proposition concerning the singularity of physical actuality. It does not presently identify that proposition with a unique, physically preferred global foliation of spacetime.

The phrase ``one current physical configuration'' must consequently be understood with care. It does not specify a universe-wide arrangement of distant events selected by an undisclosed master time coordinate. The model has not supplied such a mathematical construction.

What the proposition rejects is the necessity of treating past, present, and future as independently retained collections of completed physical configurations.

It also rejects the conclusion that delayed information or different clock readings necessarily imply multiple separately existing physical presents.

These are interpretive claims. They are not established results of special relativity, and their consistency with a complete relativistic ontology remains a theoretical issue requiring further analysis.

For purposes of this paper, the singular-present premise is retained without adding an operationally preferred universal clock.

\section{Proper Time Is Local}

The claim that physical actuality is singular must not be confused with a claim that every physical system has accumulated the same amount of elapsed time.

In relativity, proper time is associated with a timelike worldline. Using the metric signature $(-+++)$,

\begin{equation}
c^2d\tau^2=-g_{\mu\nu}dx^\mu dx^\nu.
\end{equation}

In flat spacetime, relative to a chosen inertial coordinate system, the expression becomes

\begin{equation}
d\tau=dt\sqrt{1-\frac{v^2}{c^2}}.
\end{equation}

Different timelike worldlines connecting common events may possess different proper-time lengths. If Bob and Alice depart together, follow different worldlines, and eventually reunite, the proper times recorded by their clocks may differ.

The difference is physically real. It is not merely a difference in how their experiences are described.

Yet their eventual encounter does not require their clocks to display identical accumulated intervals. The laws governing their reunion operate through their local physical conditions and interactions.

The same distinction applies within a complex physical system. The material constituents of Bob's body have different physical histories. Their organization into a coherent body does not require every constituent to have accumulated precisely the same proper time.

TLM interprets proper time as a locally accrued quantity associated with lawful physical history rather than a universally flowing temporal substance.

The singularity of physical actuality and the locality of elapsed time therefore concern different questions.

The first asks what is physically actual. The second asks what interval is accumulated along a particular timelike history.

The third relevant question is what an observer can know about distant physical conditions. That question is constrained by the propagation of information and the causal geometry of spacetime.

These three concepts must not be merged into a universal timekeeping mechanism.

\section{Space as a Registered Geometric Manifestation}

The statement ``space is not a box; it is a relation'' can now be given a more precise meaning.

The relevant relation is not merely the coordinate distance between two objects. It concerns the lawful physical relationships associated with the events under consideration. Those relationships have geometric properties that can be expressed, measured, and reconstructed through physical registration.

The governing laws are not themselves the spatial geometry. They determine the admissible structure and behavior of physical relationships. The resulting geometric relationships are expressed through actual interactions, and aspects of them become available through registration.

Space is therefore interpreted as the geometric manifestation of lawful physical relationships, not as an independent container preceding them.

The word \emph{manifestation} does not reduce geometry to subjective experience. A relationship can have objective physical consequences whether or not a human observer notices them. Tidal effects, gravitational lensing, signal timing, and detector responses are governed by lawful physical relationships that constrain what any competent measurement can reveal.

The observation-dependent aspect concerns the information available through a particular registration pathway, not the existence of the governing law.

There is also a technical limitation to the expression ``distance between events.'' In relativity, spatial distance is not generally an invariant quantity assigned uniquely to every pair of events.

In a local inertial coordinate system with signature $(-+++)$,

\begin{equation}
ds^2=-c^2dt^2+dx^2+dy^2+dz^2.
\end{equation}

The invariant interval distinguishes timelike, null, and spacelike separations. A spatial distance assigned to separated events depends on the relevant reference frame, measurement procedure, and geometric assumptions.

This distinction is not a difficulty for the relational interpretation. It shows why space should be defined through lawful geometric relationships rather than through an assumed universal grid of independent positions.

A practical example is radar ranging. Under suitable conditions, an observer sends a signal to a reflecting system and measures the proper-time interval between emission and reception. A radar-distance assignment may then be defined by

\begin{equation}
d_{\mathrm{radar}}=\frac{c\,\Delta\tau}{2},
\end{equation}

where $\Delta\tau$ is the observer's measured round-trip proper-time interval.

The assigned distance is meaningful, but it arises from a specified measurement relationship. It does not require a separately existing spatial ruler extending invisibly between the observer and the reflector.

More complicated geometries require more sophisticated measurement conventions and metric descriptions. The essential point remains that spatial quantities are established through lawful physical relationships rather than through access to an independent spatial container.

The registered manifestation of distance is one application of the governing lawful structure.

It is not the structure itself.

\section{General Relativity and the Physical Content of Geometry}

A registration-based interpretation must preserve the full physical content of general relativity.

In general relativity, the spacetime metric is dynamical. Its behavior is governed by the Einstein field equations,

\begin{equation}
G_{\mu\nu}+\Lambda g_{\mu\nu}
=\frac{8\pi G}{c^4}T_{\mu\nu}.
\end{equation}

The geometry is not merely a coordinate convention chosen by an observer. Curvature has measurable physical consequences, including tidal effects, gravitational redshift, gravitational lensing, and gravitational-wave phenomena.

Nor is the geometry reducible to pairwise distances between material bodies. Vacuum regions can possess nontrivial curvature, and gravitational waves represent propagating gravitational degrees of freedom.

The TLM interpretation must therefore distinguish the rejection of an independent spatial \emph{container} from the rejection of objectively consequential geometric structure.

The former is an ontological proposal. The latter would contradict established physical observations.

In the present account, the governing lawful structure includes the dynamical relationships expressed by general relativity. The metric describes the geometric behavior of physical systems and gravitational fields. Physical interactions and registrations reveal consequences of that geometry.

The model does not claim that a human observer creates curvature by measuring it. Rather, a measurement provides one lawful physical pathway through which the consequences of curvature become available.

This interpretation also applies to cosmology.

In a homogeneous and isotropic cosmological model, the spacetime metric may be expressed in Friedmann--Lema\^{i}tre--Robertson--Walker form:

\begin{equation}
ds^2=-c^2dt^2+
a^2(t)\left[
\frac{dr^2}{1-kr^2}
+r^2d\Omega^2
\right].
\end{equation}

Here $a(t)$ is the scale factor and $k$ specifies the spatial-curvature convention.

The expansion represented by the evolving scale factor does not require the universe to expand into a larger external spatial box. It describes changes in the metric relationships of the cosmological model.

That conclusion belongs to standard relativistic cosmology. TLM does not claim it as a new prediction.

Its additional interpretive claim is that the laws governing such metric relationships are not themselves spatial constituents residing inside the geometry they govern.

The distinction between law and geometric manifestation is therefore maintained even when the geometry is curved, dynamical, or cosmological.

\section{Motion as Lawful Relational History}

The meteor example also suggests a reinterpretation of motion.

A trajectory is a useful mathematical representation of the relationships associated with a physical system across an ordered history. It allows positions, velocities, accelerations, and encounters to be calculated within a specified spacetime description.

The success of those calculations does not, by itself, establish an independent traveler-substance passing through a completed spatial container.

Bob may describe Alice as approaching him. Alice may use coordinates in which she remains at rest and Bob approaches her. Their relative motion has measurable consequences, but neither possesses an absolute spatial velocity merely by virtue of the coordinate description.

The relational interpretation identifies the observable content of motion with lawful changes in physical relationships and the consequences registered through interactions.

This does not make physical motion an illusion.

A spacecraft can expend propellant, undergo acceleration, exchange signals, and collide or dock with another system. A meteor can deposit energy in the atmosphere and strike the ground. These are real physical consequences governed by law.

The issue is whether their description requires an additional independent spatial container and an ontologically separate substance moving through it.

TLM proposes that it does not.

This argument is conditional. The frame dependence of spatial velocity does not mathematically prove the nonexistence of independently persisting physical substances or of an independent spacetime ontology. The stronger conclusion depends on the model's premise that physical actuality is sufficiently described through lawful physical organization and interaction.

Acceleration deserves particular attention because it is not merely an arbitrary coordinate assignment.

The four-acceleration of a timelike worldline is

\begin{equation}
a^\mu=\frac{Du^\mu}{d\tau},
\end{equation}

where $u^\mu$ is four-velocity and $\tau$ is proper time.

Proper acceleration can be measured locally by an accelerometer. An observer supported on a planetary surface and an observer in free fall may record different proper accelerations, even when their spatial coordinates are represented in convenient adapted systems.

The relational account preserves these measurable facts.

It does not eliminate motion, acceleration, or the objective consequences of changing physical relationships. It rejects the requirement that these consequences be interpreted as the displacement of independent substances through a preexisting spatial box.

\section{Quantum Eligibility and Registered Outcomes}

Although the meteor is primarily a classical-scale example, the proposed lawful architecture extends to quantum phenomena.

Quantum theory describes the admissible outcomes of physical measurements through states, operators, and probability rules. For a normalized state $\lvert\psi\rangle$ and a projective measurement represented by projectors $P_i$, the Born rule assigns

\begin{equation}
p_i=\langle\psi\lvert P_i\rvert\psi\rangle.
\end{equation}

The equation provides the probability of outcome $i$ under the specified measurement conditions.

Within TLM, the mathematical structure governing admissible outcomes is interpreted as part of the timeless lawful architecture. The physical registration of a particular outcome belongs to the actual interaction.

The distinction does not require the claim that a quantum system has no physical existence between measurements. Nor does it require an observer's consciousness to select outcomes.

It also does not, by itself, solve the quantum measurement problem. Standard quantum theory and its interpretations differ on how the relationship between a state, a measurement interaction, and a definite recorded outcome should be understood.

The present paper does not replace that debate with an unsupported assertion that registration explains every aspect of quantum measurement.

Its more limited claim is that lawful eligibility and actual physical registration are distinguishable concepts, and that the distinction fits the broader separation between governing law and observable consequence.

This connection is important because it prevents the proposed ontology from being confined to human perceptions of large objects.

The same general architecture is intended to govern atomic interactions, detector outcomes, macroscopic physical organizations, and their registered spatial relationships.

The meteor is one use case of a lawful universe, not a special category of reality created by human observation.

\section{Why the Registration Interpretation Is Not Solipsism}

If spatial experience is reconstructed from present records, it may be tempting to conclude that space exists only in the mind of the observer.

That conclusion does not follow from the proposed model.

A detector is a physical system. Its recorded conditions are governed by the same laws as the events producing its signals. Independent detectors can agree because their records are constrained by common physical relationships, not because their observers share an arbitrary convention.

Bob may see the meteor crossing the sky, while a camera records it and a radar instrument measures related physical conditions. Their reports need not be identical, because their locations, instruments, and causal histories differ. Nevertheless, the reports may support a compatible reconstruction.

That agreement matters.

It distinguishes a lawfully constrained physical account from a purely private experience. It also establishes that the model cannot treat all spatial appearances as arbitrary products of memory.

The physical relationships must be sufficient to explain the observed consistency among different registration pathways.

The account consequently distinguishes a geometric relationship from a particular observer's representation of it.

The representation may be incomplete, delayed, or distorted. The lawfully governed relationships constrain which representations are possible and how competing descriptions can be compared.

The rejection of an independent box is therefore not the rejection of an objective physical world.

It is the rejection of an additional container ontology that the proposed model considers unnecessary.

\section{What the Model Establishes and What It Does Not}

The present argument has a defined logical status.

Its starting point is an ontological proposal: the totality of timeless governing laws is fundamental, physical actuality is singular, and spatial description arises as a geometric manifestation of lawfully governed relationships expressed through registration.

Several ingredients of the discussion are established physics. Special relativity describes frame-dependent spatial and temporal coordinates and invariant causal structure. General relativity supplies a dynamical account of spacetime geometry. Quantum mechanics supplies rules governing quantum states, dynamics, and measurement probabilities. Physical records are produced through interactions and interpreted through lawful inference.

Other claims are specific to the TLM interpretation. The ontological primacy of timeless governing laws is not derived from the equations quoted here. The singular universal present is not established by relativity. The absence of independently existing past configurations is not a theorem of quantum mechanics or general relativity.

The claim that space is a registered geometric manifestation also requires a distinction between explanatory interpretation and mathematical derivation.

It is possible to show that observers reconstruct spatial relationships from registered physical consequences. That fact does not prove that no independent spacetime ontology exists. A theory treating spacetime as independently real may accept the same measurement account.

TLM instead argues that an independent spatial container is unnecessary once lawful physical relationships and their geometric consequences are taken as sufficient.

The case is therefore one of ontological economy rather than a demonstrated impossibility theorem.

An additional challenge concerns the relationship between timeless law and actual physical history. Identifying laws as timeless does not by itself explain why a particular admissible outcome becomes actual, or how an ordered succession of physical registrations is obtained. The model must ultimately articulate that connection without secretly assuming the independent temporal container it seeks to dispense with.

A further challenge concerns geometry. A fully developed physical theory would need to explain how the relevant metric structure, gravitational dynamics, causal relationships, and observed quantum behavior follow from a sufficiently specified lawful foundation.

The present paper does not supply that derivation.

It provides a coherent statement of the proposed ontology, preserves the established quantitative descriptions, and identifies the conceptual questions a more complete theory would need to answer.

That limitation should not be concealed. A successful interpretation can clarify assumptions without yet producing new experimental predictions, but its scientific status must remain distinguishable from that of a derived and independently tested physical theory.

\section{Discussion: A Universe Without the Box}

The ordinary description of the meteor is simple: an object moved through space, Bob watched it, and Bob remembers the journey.

The TLM interpretation asks what lies beneath that simplicity.

The governing laws relevant to the meteor and Bob are not located in a spatial box. They constrain the admissible physical relationships and interactions of the situation. The meteor undergoes lawful interactions with its environment. Physical consequences propagate according to relativistic causal structure. Bob's visual system undergoes further lawful interactions. The resulting changes may contribute to present records from which an ordered history is reconstructed.

That history can be described as a meteor moving through space.

The description is useful, objective in its empirical constraints, and capable of quantitative refinement. What it does not independently establish is a completed spatial container through which an independently existing traveler-substance must have passed.

The same account applies to a mechanical detector. A camera records a sequence, and its present memory can be inspected. The stored frames support a reconstruction of movement. The apparatus does not need consciousness, and its records do not preserve the earlier events as independently continuing physical occurrences.

This reveals the connection between the model's treatment of space and its treatment of the past.

Space is not required as an independent container for lawful relationships. The past is not required as an independently retained collection of completed scenes for lawful records to exist.

In both cases, the ordinary picture adds a structure whose independent existence has not been established by the success of the corresponding mathematical or observational description.

The governing laws, actual physical relationships, interactions, and their present consequences remain.

Nor does the proposed universal present eliminate local proper-time differences or relativistic information delays. Physical systems possess distinct timelike histories, and registration remains constrained by causal structure. The singularity of actuality is an ontological proposition, not an invitation to introduce a master clock.

The resulting picture is not one in which the universe waits for Bob to look at it.

It is one in which Bob's observations are among the countless lawful physical consequences through which aspects of reality become registered.

The observer is not the author of physical law. The observer is one of its physical expressions.

\section{Conclusion}

Space is not required to be a completed container in which independently existing substances occupy positions and move.

The Timeless Light Model proposes that the fundamental architecture of physical reality is the totality of timeless governing laws. Those laws are not themselves constituents of space or time. They govern the admissible relationships, interactions, and outcomes of physical actuality.

A particular event brings particular lawful relationships under address. The physical consequences associated with those relationships become expressed through interactions and causally constrained registration. Some consequences are preserved as present physical records. Inspection of those records permits the reconstruction of spatial extension, movement, duration, and physical history.

Bob's experience of a meteor crossing the sky is one example. His visual system undergoes a series of lawful interactions whose preserved consequences support the description of a meteor moving through space. That description has objective physical content. It does not independently establish the existence of the spatial container imagined behind it.

The model further proposes one universal physical actuality, while distinguishing that ontological premise from a universal clock or a preferred relativistic simultaneity coordinate. Proper time is local to timelike histories. Information is available through causal registration, and delayed registration does not imply delayed existence.

Earlier events need not remain independently in existence as completed past configurations in order for their consequences to persist in present records.

This interpretation leaves established relativistic and quantum predictions intact. It is not offered as a completed mathematical derivation of spacetime geometry or as an experimental demonstration that all competing ontologies are false.

Its central contribution is the distinction between what governs physical reality and what appears through physical registration.

The timeless lawful structure is the model. The manifestation of distance, space, motion, and history is an application of that model.

The laws do not occupy the space they govern.

The physical consequences are real. Their registered relationships are measurable. Their histories can be reconstructed.

The box is not required.

\begin{thebibliography}{99}

\bibitem{LeibnizClarke}
G.~W.~Leibniz and S.~Clarke.
\newblock \emph{The Leibniz--Clarke Correspondence}.
\newblock Correspondence of 1715--1716.
\newblock H.~G.~Alexander, editor, Manchester University Press (1956).

\bibitem{Einstein1905}
A.~Einstein.
\newblock Zur Elektrodynamik bewegter K{\"o}rper.
\newblock \emph{Annalen der Physik},
\textbf{17}, 891--921 (1905).

\bibitem{EinsteinRelativity}
A.~Einstein.
\newblock \emph{Relativity: The Special and the General Theory}.
\newblock Note to the Fifteenth Edition, dated June 9, 1952,
and Appendix V, ``Relativity and the Problem of Space.''

\bibitem{Rovelli2004}
C.~Rovelli.
\newblock Space Does Not Exist.
\newblock In \emph{2004: What's Your Law?},
Edge Foundation (2004).
\newblock \url{https://www.edge.org/response-detail/11779}.

\bibitem{Rovelli1996}
C.~Rovelli.
\newblock Relational Quantum Mechanics.
\newblock \emph{International Journal of Theoretical Physics},
\textbf{35}, 1637--1678 (1996).

\bibitem{Wald}
R.~M.~Wald.
\newblock \emph{General Relativity}.
\newblock University of Chicago Press (1984).

\bibitem{TaylorWheeler}
E.~F.~Taylor and J.~A.~Wheeler.
\newblock \emph{Spacetime Physics}.
\newblock W.~H.~Freeman, 2nd edition (1992).

\bibitem{Zurek2003}
W.~H.~Zurek.
\newblock Decoherence, einselection,
and the quantum origins of the classical.
\newblock \emph{Reviews of Modern Physics},
\textbf{75}, 715--775 (2003).

\end{thebibliography}

\end{document}
